\chapter{Conclusion and Future Work}
\label{ch:conclusion}

\section{What the study found}
This thesis began with a reasonable expectation: questions that need
different amounts of evidence should benefit from different retrieval
granularities. The experiments support the premise that the best size
varies across questions, but they do not establish that the evaluated
routers can exploit that variation better than a strong fixed setting.

On the frozen validation set, fixed 40-token retrieval reaches mean
joined evidence-token F1 of 0.3159. The offline best-of-five oracle
reaches 0.3821, leaving a meaningful gap. The strongest learned
question-level method, expected-regret routing, reaches 0.3075.
Local tree-level routing reaches 0.3053, compared with 0.3075 for
its matched fixed-40 tree baseline. Adding a frozen Qwen
context-need feature does not demonstrate an improvement.

The sequence separates the promise of adaptation from the information
and supervision needed to turn it into a reliable decision.

\section{Answers to the research questions}
\subsection{RQ1: Is there headroom beyond a fixed size?}
Yes, within the five saved retrieval actions. The offline oracle
improves mean F1 by 0.0662 over fixed 40, and its selected sizes
cover the full candidate set. A single size is therefore not
optimal for every question. However, oracle headroom alone does
not show that the correct decision is predictable without access
to gold evidence. It establishes opportunity, not deployability.

\subsection{RQ2: Is the question alone sufficient?}
The evaluated question-only methods do not beat fixed 40.
The embedding router reaches 0.2853 retrieval F1, while the
strongest retrieval result among the reported Qwen question-only
variants is 0.2865 for uniform alias training. Fine-tuning and
classification-head changes improve some label metrics without
consistently improving retrieval.

The clearest counterexample is numeric SFT: label accuracy
increases from 0.0400 to 0.4318 relative to frozen prompting,
but retrieval F1 decreases. This shows that learning the
evidence-length category is not an adequate substitute for
learning a useful retrieval action in this protocol.

\subsection{RQ3: Does similarity structure help?}
Similarity-tree features provide additional observable information.
The weighted boosted-tree setup improves mean retrieval F1 and macro
F1 over the tested uniformly weighted linear tree model, but does not
isolate a benefit from nonlinearity alone. Original fusion looks promising,
but its in-sample base features make the training-fold estimate
optimistic. Corrected OOF fusion reaches 0.2783 retrieval F1.

The evidence therefore supports investigating combined language
and document signals, but not claiming that hierarchy or fusion
alone solves granularity selection. The corrected training
procedure is a methodological contribution even though it lowers
the apparent result.

\subsection{RQ4: Does retrieval-utility supervision help?}
Yes, relative to the preceding learned global methods in the
recorded development comparisons. Expected-regret learning
reaches 0.3075 F1, with positive paired differences over
Phase 2D, Phase 3B, and both fusion versions. The change of
objective also produces a markedly different prediction
distribution, concentrating on 20 and 40 rather than 80 and 160.

It still falls short of fixed 40 by 0.00835 F1. The result
supports objective alignment as a productive direction, not
a final superiority claim for the learned router.

\subsection{RQ5: Does local adaptation add value?}
The local protocol allows different regions to contribute
different chunk sizes, but the trained router does not demonstrate
an advantage over the same-tree fixed-40 baseline. Its slight
recall gain comes with lower precision. Returning all chunks
at each predicted level raises recall much further while
reducing F1 and substantially increasing the output budget.

The frozen Qwen context feature also adds no demonstrated gain.
Constrained scoring successfully removes invalid outputs, but
almost every question receives the category medium. A reliable
interface is necessary; it does not guarantee a discriminative
feature.

\section{Implications for adaptive retrieval}
The most useful design principle is to define the retrieval action
before its label. If the action returns five chunks, the target
should reflect their utility. Total annotation length is an
intuitive property, but it is not the action itself.

A second principle is to account for information cost. A question-only
router needs less preparation than one that inspects every level
of a paper. Additional features must earn their cost through a
measurable quality or efficiency benefit.

A third principle is to retain strong controls as the research
becomes more sophisticated. The global fixed-size baseline,
the matched local baseline, and the OOF reconstruction each
prevent a different misleading comparison. They are part of
the scientific result, not merely implementation details.

\section{Future work}
\subsection{Independent evaluation and repeated training}
The immediate priority is an untouched evaluation split with
a frozen protocol. The action set, target, feature extraction,
hyperparameters, checkpoint rule, and primary metric should
be fixed before inspecting its outcomes. Repeated training
seeds would then distinguish stable gains from a favourable
single fit.

If further model selection is required, it should occur inside
paper-grouped training procedures, with nested construction
for stacked models where appropriate. The external evaluation
should be used once for the specified comparison. Additional
datasets would test whether the observed preference for an
intermediate size is specific to QASPER and this metric.

\subsection{Utility-aware local decisions}
The most direct modelling extension is to bring the Phase 4
objective into the local setting. Instead of predicting local
evidence length, a tree-level learner could estimate the
incremental retrieval utility of each candidate chunk. This
would also expose an interaction that the current local
classifier ignores: a useful chunk may become redundant after
another tree has already contributed similar evidence.

A set-level decision under an explicit token budget would
address that interaction more naturally. Candidate selection
could stop when further text adds little expected evidence
coverage, rather than always returning five independent
decisions. This is a proposed experiment, not an implemented
result of the repository.

\subsection{Calibration and selective adaptation}
Most questions in Phase 4 are assigned 20 or 40, while some
large-size oracle opportunities are missed. A selective router
could use fixed 40 as a default and deviate only when the
predicted utility difference is sufficiently large and
well-calibrated. The threshold would need training-only
selection, and performance should be reported as a function
of both deviation frequency and retrieval cost.

The question is whether the system can identify when changing
a strong default is worthwhile. The current results motivate
that experiment without answering it.

\subsection{Representations, budgets, and generation}
Sentence-aware or proposition-based units could reduce the
boundary failures of short token windows. Contextual chunk
embeddings could preserve neighbouring information without
always returning longer passages. Both changes should be
evaluated with matched candidate and token budgets, since
otherwise a representation effect can be confused with an
increase in returned text.

A complete RAG evaluation should then test whether changes in
retrieved evidence improve generated answers. It should hold
the generator, prompt, context ordering, and generation
settings constant, and measure answer correctness, evidence
support, and citation quality alongside retrieval F1. Questions
requiring tables or figures deserve separate analysis rather
than exclusion from every stage.

Finally, a deployment-oriented study should measure indexing
cost, feature-extraction latency, model memory, retrieval
latency, and total supplied tokens under comparable hardware
conditions. An adaptive method can be useful by saving
resources at similar quality, even if its absolute F1 does
not exceed the best fixed baseline. That trade-off was not
established by the present experiments.

\section{Closing remarks}
Adaptation must earn its place against a strong fixed baseline.
Here, utility-aware supervision and cleaner evaluation made progress,
but did not establish superiority. The next step is to test that
direction on untouched data under a fixed protocol.
